%%=============================================================================
%% Inleiding
%%=============================================================================

\chapter{\IfLanguageName{dutch}{Inleiding}{Introduction}}%
\label{ch:inleiding}

Een bezoek aan een kerkhof is voor veel mensen een emotioneel moment van herdenking en rust. Nabestaanden willen in alle vrede het graf van een dierbare kunnen terugvinden, zonder bijkomende stress of verwarring. In de praktijk blijkt dit echter niet altijd vanzelfsprekend. Grote of complex ingerichte begraafplaatsen, onduidelijke nummering, ontbrekende of moeilijk leesbare grafstenen en weersomstandigheden kunnen dit een stuk moeilijker maken. Hierdoor wordt dit emotioneel bezoek meer een frustrerende puzzel. Tegelijkertijd is digitalisatie steeds meer aanwezig in het dagelijks leven. Navigatietoepassingen op smartphones zijn uitgegroeid tot een onmisbaar hulpmiddel in steden, gebouwen en natuurgebieden. Ook binnen lokale besturen worden begraafplaatsen gedigitaliseerd, dit via databanken en geoloketten. Toch zijn deze systemen niet altijd rechtstreeks toegankelijk voor bezoekers, of beperken ze zich tot statische kaarten zonder interactieve of contextgevoelige ondersteuning. Hierdoor blijft de kloof bestaan tussen beschikbare digitale data en de effectieve gebruikservaring op het terrein.

\section{\IfLanguageName{dutch}{Probleemstelling}{Problem Statement}}%
\label{sec:probleemstelling}

Hoewel er reeds digitale kaarten en webtoepassingen bestaan voor bepaalde begraafplaatsen, ervaren bezoekers nog steeds problemen bij het efficiënt navigeren naar een specifiek graf. Klassieke kaartweergaven vereisen dat gebruikers hun aandacht afwisselend richten tussen het scherm en de fysieke omgeving. Dit kan verwarrend zijn, zeker voor gebruikers die minder ervaring hebben met zo’n digitale applicatie. Augmented Reality (AR) kan hier mogelijk een meer intuïtieve oplossing voor bieden. Door digitale beelden en informatie rechtstreeks te projecteren op de werkelijke omgeving via de camera, kan de gebruiker visuele aanwijzingen volgen zonder voortdurend te moeten schakelen tussen kaart en omgeving. Toch is het onduidelijk in welke mate AR effectief bijdraagt aan een rustigere, efficiëntere en een meer gebruiksvriendelijkere navigatie-ervaring op een kerkhof. Daarnaast moeten we ons vragen stellen rond privacy, datagebruik en ethische grenzen als we over digitalisatie van begraafplaatsgegevens spreken.

\section{\IfLanguageName{dutch}{Onderzoeksdoelstelling}{Research objective}}%
\label{sec:onderzoeksdoelstelling}

In deze bachelorproef wordt er onderzocht hoe digitalisatie, in combinatie met Augmented Reality, kan worden ingezet om de efficiëntie en gebruikservaring van kerkhofnavigatie te verbeteren. Het onderzoek focust zich op de ontwikkeling en evaluatie van een Proof of Concept: een mobiele applicatie voor iOS die een digitale kaart combineert met een AR-overlay voor routebegeleiding naar een specifiek graf. Ook voor Android toestellen worden de mogelijkheden onderzocht, maar de focus van de proof of concept ligt op mobiele iOS toestellen. De studie beperkt zich tot mobiele toestellen van zowel iOS en Android die gebruik kunnen maken van AR-functionaliteiten voor positionele tracking en visuele navigatie. Er wordt geen volledige integratie met gemeentelijke administratiesystemen gerealiseerd, maar er wordt gewerkt met bestaande of gesimuleerde data om de technische en gebruiksgerichte haalbaarheid te onderzoeken. Daarnaast wordt expliciet aandacht besteed aan privacyaspecten en aan de noden van verschillende leeftijdsgroepen.


\section{\IfLanguageName{dutch}{Onderzoeksvraag}{Research question}}%
\label{sec:onderzoeksvraag}

De centrale onderzoeksvraag:

\begin{center}
    \vfill
    \textit{Op welke manier kan digitalisatie, in combinatie met Augmented Reality, worden ingezet om de efficiëntie van het navigeren op een kerkhof te verbeteren?}
    \vfill
\end{center}

\vspace{0.5cm}

Aanvullend worden er een aantal deelvragen gesteld die het onderzoek verder ondersteunen:

\begin{itemize}
    \item Welke moeilijkheden ervaren bezoekers bij het navigeren op kerkhoven?
    \item Welke beperkingen hebben bestaande navigatietoepassingen?
    \item Welke behoeften hebben bezoekers van verschillende leeftijden?
    \item Welke emotionele en contextuele factoren spelen een rol?
    \item Welke privacygevoelige data is hierbij betrokken?
    \item Hoe kan Augmented Reality ondersteuning bieden bij kerkhofnavigatie?
    \item Hoe kan AR gecombineerd worden met digitale kaarten of geoloketten?
    \item Welke ontwerprichtlijnen zijn nodig om rust en sereniteit te bewaren?
    \item Welke data kan veilig gebruikt worden?
    \item Hoe kan deze Proof of Concept getest en geëvalueerd worden?
\end{itemize}


\subsection{Relevantie van het onderzoek}
Dit onderzoek is maatschappelijk relevant omdat het zich afspeelt binnen een concrete, emotioneel gevoelige context waarin gebruiksvriendelijkheid en respect centraal staan. Door technologie op een efficiënte manier in te zetten, kan een bezoek aan een begraafplaats op een betere en minder stressvolle manier verlopen. Indien het Proof of Concept positieve resultaten oplevert, kan dit een eerste stap vormen naar een bredere implementatie binnen lokale besturen en bijdragen aan de verdere digitalisering van publieke dienstverlening.

\section{\IfLanguageName{dutch}{Opzet van deze bachelorproef}{Structure of this bachelor thesis}}%
\label{sec:opzet-bachelorproef}

% Het is gebruikelijk aan het einde van de inleiding een overzicht te
% geven van de opbouw van de rest van de tekst. Deze sectie bevat al een aanzet
% die je kan aanvullen/aanpassen in functie van je eigen tekst.

De rest van deze bachelorproef is als volgt opgebouwd:

In Hoofdstuk~\ref{ch:stand-van-zaken} wordt een overzicht gegeven van de stand van zaken binnen het onderzoeksdomein, op basis van een literatuurstudie.

In Hoofdstuk~\ref{ch:methodologie} wordt de methodologie toegelicht en worden de gebruikte onderzoekstechnieken besproken om een antwoord te kunnen formuleren op de onderzoeksvragen.

In Hoofdstuk~\ref{ch:voorbereiding-analyse} wordt de voorbereidingsfase
beschreven: de risicoanalyse, het interview met de gemeente Sint-Gillis-Waas
over het systeem WebBGP, en de terreinanalyse van de testlocatie.

In Hoofdstuk~\ref{ch:ontwerp} wordt het ontwerp van de applicatie beschreven:
de gelaagde softwarearchitectuur, de gebruikersinterface en -flow, en de
AR-navigatielogica.

In Hoofdstuk~\ref{ch:juridische-analyse} wordt een juridische en ethische
analyse uitgevoerd van het gebruik van persoonsgegevens binnen de applicatie,
met bijzondere aandacht voor de GDPR, het rijksregisternummer en de rolverdeling
tussen de betrokken partijen.

In Hoofdstuk~\ref{ch:implementatie} wordt de technische uitwerking van de
applicatie beschreven, inclusief de AR-component, de GPS-naar-AR coördinatenomzetting,
de fallback-mechanismen en de platformonafhankelijkheid.

In Hoofdstuk~\ref{ch:testen} worden de testopzet en de resultaten van de
technische evaluatie op het kerkhof van Sint-Gillis-Waas beschreven.

In Hoofdstuk~\ref{ch:analyse-evaluatie} worden de testresultaten geanalyseerd,
de onderzoeksvragen beantwoord en aanbevelingen geformuleerd voor verdere
doorontwikkeling.

In Hoofdstuk~\ref{ch:conclusie}, tenslotte, wordt de conclusie gegeven en een antwoord geformuleerd op de onderzoeksvragen. Daarbij wordt ook een aanzet gegeven voor toekomstig onderzoek binnen dit domein.